\documentclass[11pt,a4paper]{article}

\usepackage[utf8]{inputenc}
\usepackage[T1]{fontenc}
\usepackage[spanish]{babel}
\usepackage{lmodern}
\usepackage{geometry}
\usepackage{booktabs}
\usepackage{tabularx}
\usepackage{longtable}
\usepackage{array}
\usepackage{enumitem}
\usepackage[hidelinks]{hyperref}

\geometry{margin=2.2cm}
\setlength{\parindent}{0pt}
\setlength{\parskip}{0.55em}
\renewcommand{\arraystretch}{1.16}
\newcolumntype{Y}{>{\raggedright\arraybackslash}X}
\newcommand{\alta}{\textbf{Alta}}
\newcommand{\media}{Media}
\newcommand{\opcional}{Opcional}

\title{Planificación corregida de próximos lotes\\de mosto natural, 2026}
\author{Campaña de discriminación estructural}
\date{9 de septiembre de 2026}

\begin{document}
\maketitle

\begin{abstract}
Se propone una campaña bloqueada de dos lotes, cada uno formado por tres
fermentadores simultáneos. En ambos lotes aparecen el control sin nutrición y el
tratamiento con nutrición conocida y trazable. La campaña reúne seis unidades
experimentales: \(n=3\) controles y \(n=3\) nutridos. El lote es el bloque
experimental, el tratamiento es el factor científico y cada fermentador es una
unidad experimental. La asignación física se contrabalancea para que F1, F2 y F3
cambien de tratamiento entre lotes.
\end{abstract}

\section{Objetivo}

El objetivo es generar un contraste causal que separe terminación basal,
respuesta nutricional rápida, efecto nutricional sostenido y dinámica lenta de
terminación. Los datos también deben mostrar si el CO$_2$ online contiene
información suficiente para inferir posteriormente \(X\) y \(N\), usando
Oculyze, biomasa y química offline como verdad de referencia.

\section{Estado actual del modelo}

El modelo central contiene \(X,X_d,N,G,F,E,\mathrm{Gly}\); la temperatura es un
driver exógeno. \(N\) limita el crecimiento, pero no todas las rutas de consumo
de azúcares. La muerte celular queda prácticamente inactiva en el rango térmico
relevante, pese a la señal de \(X_d\) observada. La capa de CO$_2$ combina
activación, pool disuelto, liberación continua y una ganancia empírica de matriz.

La respuesta nutricional vigente mezcla un cambio post-pulso lento con la
terminación. Los diagnósticos recientes sugieren un transitorio rápido separado,
pero todavía no distinguen con datos independientes respuesta rápida, efecto
sostenido y terminación basal. El onset depende además de química offline futura
y no es completamente causal para un estimador online.

\section{Por qué se necesitan nuevos experimentos}

LAB013--015, cerca de \(16\,^{\circ}\mathrm{C}\), aportan química y buena
cobertura Oculyze, pero CO$_2$ gaseoso menos confiable. LAB016--018, cerca de
\(20\,^{\circ}\mathrm{C}\), aportan CO$_2$ confiable, aunque casi sin ground
truth temporal de \(X,X_d,N,G,F,E,\mathrm{Gly}\). Sus pulsos ocurrieron cerca de
56.93 h, pero la formulación y dosis exactas no están confirmadas.

Se requieren por ello fermentaciones donde tratamiento, tiempo, dosis, química,
biomasa y CO$_2$ estén sincronizados. LAB016--018 se mantienen como holdout
histórico y no se convierten en training durante esta tarea.

\section{Hipótesis a discriminar}

\begin{enumerate}[leftmargin=*]
  \item El control identifica onset y terminación basal sin el artefacto de un
  evento nutricional.
  \item Una divergencia rápida tratamiento--control post-pulso evidencia
  respuesta nutricional rápida; su persistencia evidencia efecto sostenido.
  \item La reconvergencia apoya un efecto principalmente transitorio.
  \item Fructosa residual y CO$_2$ de cola permiten distinguir metabolismo
  upstream de memoria exclusivamente gaseosa.
  \item Un aumento medido de \(X_d\) con predicción \(X_d\simeq0\) exige revisar
  la dinámica de muerte.
  \item Cambios de CO$_2$ incompatibles con \(X,N,G,F,E\) y CO$_2$ disuelto
  señalan la capa de observación, transferencia o \texttt{matrix\_gain}.
\end{enumerate}

\section{Diseño experimental elegido}

La campaña tiene la siguiente jerarquía:
\[
\begin{array}{l}
\text{Campaña de discriminación estructural}\\
\quad\text{Lote 1: F1, F2, F3}\\
\quad\text{Lote 2: F1, F2, F3.}
\end{array}
\]

\begin{table}[ht]
\centering
\small
\begin{tabularx}{\textwidth}{@{}l l Y c l l@{}}
\toprule
Lote & Fermentador & Tratamiento & Réplica & Inoculación & Pulso/pseudo-pulso\\
\midrule
Lote 1 & F1 & Control sin nutrición & C1 & 21-09-2026 11:55 & 24-09-2026 09:40\\
Lote 1 & F2 & Control sin nutrición & C2 & 21-09-2026 12:20 & 24-09-2026 10:05\\
Lote 1 & F3 & Nutrición conocida & N1 & 21-09-2026 12:45 & 24-09-2026 10:30\\
\midrule
Lote 2 & F1 & Nutrición conocida & N2 & 05-10-2026 11:55 & 08-10-2026 09:40\\
Lote 2 & F2 & Nutrición conocida & N3 & 05-10-2026 12:20 & 08-10-2026 10:05\\
Lote 2 & F3 & Control sin nutrición & C3 & 05-10-2026 12:45 & 08-10-2026 10:30\\
\bottomrule
\end{tabularx}
\caption{Diseño bloqueado y contrabalanceado. Cada lote contiene ambos
tratamientos y cada fermentador físico cambia de condición entre lotes.}
\end{table}

El resultado global es exactamente \(n=3\) controles y \(n=3\) nutridos. Todos
operan a \(20.0\,^{\circ}\mathrm{C}\). En los controles, el pseudo-pulso es sólo
una marca temporal: no se añade producto ni líquido.

\section{Justificación del diseño bloqueado}

El tratamiento es el factor cuya consecuencia se desea estimar. El lote agrupa
condiciones comunes de mosto, preparación, fecha y operación, por lo que se trata
como bloque. El fermentador es la unidad experimental independiente que recibe
el tratamiento.

No se asignan tres controles al Lote 1 y tres nutridos al Lote 2 porque entonces
tratamiento y lote serían perfectamente confundidos: cualquier diferencia de
mosto, preparación, instrumento o fecha podría interpretarse erróneamente como
efecto nutricional. La distribución \(2C+1N\) en Lote 1 y \(1C+2N\) en Lote 2
permite buscar el contraste dentro de cada bloque y comprobar su consistencia
entre bloques. Los triplicados son globales por condición, no tres réplicas
dentro de cada lote.

\section{Asignación física y contrabalanceo}

Se fija Lote 1: F1=C, F2=C, F3=N; Lote 2: F1=N, F2=N, F3=C. Así F1, F2 y F3
reciben cada tratamiento exactamente una vez durante la campaña. Este
contrabalanceo garantiza que ningún fermentador quede siempre asociado a una
condición y es preferible, con sólo dos lotes, a una randomización irrestricta
que podría repetir accidentalmente la misma asociación. La asignación debe
cerrarse antes de iniciar y cualquier cambio debe registrarse.

\section{Justificación de temperatura}

Se seleccionan \(20\,^{\circ}\mathrm{C}\) para las seis fermentaciones. Esto
elimina temperatura como factor concurrente, permite comparación térmica con
LAB016--018 y concentra el contraste en nutrición y terminación. La condición
aproximada de \(18\,^{\circ}\mathrm{C}\) se reserva para un holdout posterior.

\section{Tratamientos y nutrición}

Los controles no reciben nutrición. Los nutridos reciben una única formulación
confirmada antes de ejecutar. Deben registrarse producto, fabricante/tipo, lote
de producto, masas de SFX y FDA, volumen inicial y actual, dosis calculada,
fecha/hora real de inicio y fin, duración, operador y observaciones. No se fija
un \(\Delta N\), porcentaje YAN ni masa final sin respaldo documental.

El criterio es fijo, observable y causal:
\[
t_{\mathrm{pulse}}=69.75\ \mathrm{h}
\]
desde el fin de inoculación de cada fermentador. Se aparta 12.82 h del evento
de LAB016--018 porque mantener aproximadamente 56--57 h desplazaría el pulso al
final de la jornada y perdería la ventana manual rápida. Brix, densidad y CO$_2$
previos documentan el estado fisiológico, pero la química futura no decide la
intervención.

\section{Variables y prioridades}

\begin{longtable}{@{}p{0.25\textwidth}p{0.13\textwidth}p{0.54\textwidth}@{}}
\toprule
Medición & Prioridad & Función científica\\
\midrule
\endhead
CO$_2$ gaseoso continuo & \alta & Señal online: onset, respuesta rápida, área y
cola; requiere blanco, calibración, reloj común e interrupciones documentadas.\\
Temperatura continua & \alta & Input real y comprobación de igualdad térmica.\\
Oculyze (\(X/X_d\)) & \alta & Crecimiento, máximo, viabilidad y muerte; réplicas
técnicas en puntos clave si son viables.\\
Biomasa independiente & \alta & Ground truth de \(X\) y control de conversión
Oculyze--masa.\\
\(G,F\) & \alta & Consumo y fructosa residual de cola.\\
YAN/componentes & \alta & Trayectoria de \(N\) y efecto efectivo del producto.\\
Etanol \(E\) & \alta & Producción metabólica y fuente de la capa CO$_2$.\\
Brix/densidad & \alta & Avance causal; no convertir automáticamente en \(G/F\)
sin función de observación validada.\\
Nutrición/volumen retirado & \alta & Inputs y balance de volumen por unidad.\\
Glicerol & \media & Estado modelado y balance secundario.\\
CO$_2$ disuelto & \media & Pocos puntos dirigidos para separar producción,
almacenamiento y liberación.\\
O$_2$ disuelto & \opcional & Variable auxiliar de la capa CO$_2$; usar sólo una
sonda validada sin comprometer el núcleo del plan.\\
\bottomrule
\end{longtable}

\section{Diseño de muestreo}

Todos los fermentadores comparten los mismos tiempos nominales. Las horas
programadas están separadas 25 min y se conserva siempre la hora real. Las
fechas de la tabla se expresan como Lote 1 / Lote 2.

\begin{longtable}{@{}r l l l l l@{}}
\toprule
\(t\) [h] & Relación & Fecha L1/L2 & F1 & F2 & F3\\
\midrule
\endhead
-1.00 & Inicio automático & 21-09 / 05-10 & 10:55 & 11:20 & 11:45\\
0.25 & Condición inicial & 21-09 / 05-10 & 12:10 & 12:35 & 13:00\\
20.25 & Crecimiento/onset & 22-09 / 06-10 & 08:10 & 08:35 & 09:00\\
44.25 & Crecimiento/pico & 23-09 / 07-10 & 08:10 & 08:35 & 09:00\\
68.25 & Pre-pulso \(-1.5\) h & 24-09 / 08-10 & 08:10 & 08:35 & 09:00\\
69.75 & Pulso/pseudo-pulso & 24-09 / 08-10 & 09:40 & 10:05 & 10:30\\
71.75 & Post \(+2\) h & 24-09 / 08-10 & 11:40 & 12:05 & 12:30\\
73.75 & Post \(+4\) h & 24-09 / 08-10 & 13:40 & 14:05 & 14:30\\
75.25 & Post \(+5.5\) h & 24-09 / 08-10 & 15:10 & 15:35 & 16:00\\
93.75 & Post \(+24\) h & 25-09 / 09-10 & 09:40 & 10:05 & 10:30\\
164.25 & Terminación/cola & 28-09 / 12-10 & 08:10 & 08:35 & 09:00\\
212.25 & Cola tardía & 30-09 / 14-10 & 08:10 & 08:35 & 09:00\\
260.25 & Final planificado & 02-10 / 16-10 & 08:10 & 08:35 & 09:00\\
\bottomrule
\end{longtable}

Hay 11 inicios manuales por fermentación, 33 por lote y 66 en toda la campaña.
Cada lote añade los eventos de nutrición correspondientes, mientras los
pseudo-pulsos no requieren intervención. El libro Excel detalla las técnicas de
cada punto y deja vacíos tiempo real y volumen.

\section{Ventana post-pulso y carga operativa}

Con tres fermentadores, se conservan pre-pulso, \(+2,+4,+5.5,+24\) h. El punto
\(+4\) h es simple; los demás son complejos. Se retira \(+1\) h manual porque
generaría solapamiento con el conjunto \(+2\) h y se mantiene CO$_2$/temperatura
continuos para capturar esa respuesta. No se programa \(+12\) h porque cae de
noche. Esta reducción afecta primero puntos manuales redundantes y no el número
de réplicas ni las anclas de \(X/X_d\), \(N\), sustratos y terminación.

Los conjuntos complejos se inician cada 25 min. El último punto complejo del
jueves comienza a las 16:00 y termina estimativamente a las 16:25, dejando cerca
de una hora para almacenamiento, limpieza y registro. El viernes, el último
comienza a las 10:30 y termina cerca de las 10:55.

\section{Biomasa y Oculyze}

Oculyze cubre condición inicial, crecimiento/onset, cercanía al máximo,
pre-pulso, \(+2,+5.5,+24\) h, terminación, cola y final. Se obtienen diez anclas
por fermentación, comparables con la cobertura de LAB013--015 y muy superiores a
las tres observaciones de LAB016--018. Se conservan concentración, viabilidad y
metadatos crudos; las réplicas técnicas no reemplazan las seis unidades
experimentales.

\section{Planificación calendario}

Lote 1 se inocula el lunes 21 de septiembre y termina de forma planificada el
viernes 2 de octubre. Lote 2 replica el patrón de días y horas desde el lunes 5
hasta el viernes 16 de octubre. La separación permite reutilizar F1--F3 y bloquear
fecha/preparación. Debe confirmarse que el intervalo 2--5 de octubre basta para
limpieza y preparación; si no, todo el Lote 2 se desplaza una semana sin deformar
su grilla.

Iniciar antes del 18 de septiembre se descarta: el pulso o su ventana rápida
caerían en la jornada reducida del 17 o en el feriado del 18.

\section{Restricciones operativas}

\begin{table}[ht]
\centering
\small
\begin{tabularx}{\textwidth}{@{}l l l Y@{}}
\toprule
Día/fecha & Presencia & Inicio manual conservador & Aplicación\\
\midrule
Lunes--jueves & 08:00--17:30 & 08:10--16:30 & Se cumple; último inicio 16:00.\\
Viernes & 08:00--12:00 & 08:10--11:00 & Se cumple; último inicio 10:30.\\
16-09-2026 & 08:00--15:30 & 08:10--14:30 & Ningún lote activo.\\
17-09-2026 & 08:00--12:00 & 08:10--11:00 & Ningún lote activo.\\
18-09-2026 & Feriado & Sin trabajo manual & Ningún lote activo.\\
\bottomrule
\end{tabularx}
\end{table}

\section{Relación con recalibración}

No se recalibra nada en esta tarea. Tras cerrar los datos, la campaña podrá
apoyar selección estructural y una recalibración justificada. Antes se debe
adoptar una activación causal, separar transitorio y terminación, unificar la
convención de CO$_2$ y confirmar la semántica/dosis de \(N\).

\section{Relación con el estimador}

El estimador futuro recibirá CO$_2$, temperatura y eventos conocidos online; no
puede usar química futura. El diseño bloqueado permite evaluar si el contraste
visible en CO$_2$ se repite dentro de ambos lotes y coincide con cambios reales
de \(X,X_d,N,G,F,E\). La falta de coherencia debe resolverse antes de confiar en
estimaciones de \(X/N\).

\section{Criterios de éxito}

\begin{itemize}[leftmargin=*]
  \item Una diferencia nutrición--control consistente dentro de ambos lotes
  apoya un efecto de tratamiento y no sólo de lote.
  \item Una divergencia rápida seguida de reconvergencia apoya un transitorio;
  persistencia a \(+24\) h/cola apoya efecto sostenido.
  \item Una cola anómala también en controles apunta a terminación basal/upstream.
  \item \(X_d\) creciente con predicción casi nula rechaza la muerte vigente.
  \item CO$_2$ sin correspondencia metabólica obliga a revisar observación,
  transferencia o ganancia.
  \item Disminuir RMSE sin separar estas alternativas no basta.
\end{itemize}

\section{Riesgos y contingencias}

\begin{itemize}[leftmargin=*]
  \item \textbf{Dosis no confirmada:} no iniciar condiciones nutridas ni asignar
  \(\Delta N\).
  \item \textbf{Carryover entre lotes:} validar limpieza de F1--F3 antes del
  Lote 2.
  \item \textbf{Ciclo complejo mayor de 25 min:} reprogramar el lote completo o
  retirar primero el punto de prioridad media, sin perder réplicas.
  \item \textbf{Fuga/reloj/sensor:} ejecutar blancos y hermeticidad; documentar
  interrupciones sin interpolarlas silenciosamente.
  \item \textbf{Volumen excesivo:} validar volumen por técnica y máximo acumulado.
  \item \textbf{Desviación horaria:} registrar hora y tiempo reales.
  \item \textbf{Final temprano/tardío:} aplicar un criterio predefinido y trasladar
  el conjunto final, manteniendo CO$_2$/T continuos.
\end{itemize}

\section{Futuro holdout a \(18\,^{\circ}\mathrm{C}\)}

Después de seleccionar y, si corresponde, recalibrar la estructura con esta
campaña, se recomienda un lote independiente cerca de \(18\,^{\circ}\mathrm{C}\)
para validar interpolación entre 16 y 20 \(^{\circ}\mathrm{C}\). Debe reservarse
como validación externa y no añadirse a los dos lotes actuales.

\section{Decisiones pendientes antes de ejecutar}

\begin{enumerate}[leftmargin=*]
  \item Confirmar producto, fabricante/tipo/lote, composición, masas SFX/FDA,
  aporte YAN y duración de aplicación.
  \item Confirmar volumen inicial, headspace, volumen por técnica y máximo
  acumulado admisible.
  \item Confirmar equivalencia, disponibilidad y limpieza de F1/F2/F3 para ambos
  lotes.
  \item Confirmar reactivos, Oculyze, biomasa y capacidad analítica para tres
  fermentadores dentro de ciclos de 25 min.
  \item Confirmar CO$_2$ disuelto y O$_2$ opcional.
  \item Cerrar blancos/calibración de caudalímetros, relojes, hermeticidad,
  frecuencia de logging y convención de unidades.
  \item Definir criterio causal de cierre mediante CO$_2$ sostenido y estabilidad
  de Brix/densidad.
\end{enumerate}

\section*{Control de consistencia}

El Excel contiene exactamente dos lotes, tres fermentadores por lote, seis
fermentaciones, \(n=3\) por tratamiento, ambas condiciones en ambos lotes y los
mismos tiempos nominales. No hay operaciones manuales el 18 de septiembre. Las
fechas se calcularon sumando cada tiempo nominal a la inoculación correspondiente.

\end{document}
